\chapter{Pendahuluan}

\section{Latar Belakang}
Layanan jaringan nirkabel (WiFi) di lingkungan perguruan tinggi umumnya dilindungi mekanisme \textit{captive portal} berbasis web agar hanya sivitas akademika yang memiliki akun resmi yang dapat mengakses internet \cite{rfc8952}. Kendati efektif dari sisi pengendalian akses, proses autentikasi manual melalui peramban menimbulkan friksi ketika pengguna berpindah gedung, sesi kedaluwarsa, atau perangkat terputus sementara dari titik akses (\textit{access point}). Pengguna harus membuka ulang halaman portal dan mengetikkan nama pengguna serta kata sandi secara berulang.

Aplikasi \textbf{Coneckin} hadir sebagai solusi Android native berbasis Kotlin dan Jetpack Compose yang memungkinkan proses masuk ke portal WiFi kampus cukup dengan satu ketukan, baik dari layar utama aplikasi maupun lewat \textit{Quick Settings Tile}. Untuk menjaga privasi, kredensial pengguna disimpan terenkripsi lokal di perangkat memakai AES-256 GCM melalui \texttt{AndroidKeyStore} tanpa pernah dikirim ke server pusat \cite{android_keystore,nist_gcm}. Sementara itu, layanan server REST API pada Node.js dan MongoDB menyediakan direktori portal kampus terverifikasi, informasi konfigurasi rilis aplikasi, dan penampung laporan gangguan anonim bersifat opsional (\textit{opt-in}).

Pada versi awal aplikasi, masih terdapat beberapa kesenjangan rekayasa perangkat lunak yang perlu diselesaikan dalam Proyek 2:
\begin{enumerate}[leftmargin=*]
  \item Operasi \textit{Update} pada penyimpanan akun lokal Android belum tersedia sehingga pengguna harus menghapus dan membuat ulang akun ketika kata sandi atau alamat portal berubah.
  \item Layanan server API sebelumnya baru menyediakan operasi baca publik (\texttt{GET /portals}, \texttt{GET /config}) dan tambah laporan (\texttt{POST /reports}), belum memiliki endpoint CRUD administratif lengkap maupun header CORS bagi konsumsi web.
  \item Pengujian otomatis belum mencakup pengukuran \textit{code coverage}, belum menyertakan pengujian instrumentasi untuk \texttt{AndroidKeyStore}, dan alur CI GitHub Actions belum menegakkan ambang kelulusan uji sebelum membangun rilis APK/AAB.
  \item Belum tersedianya halaman panduan dan direktori portal berbasis konvensi modular Croot.js pada repositori terpisah \texttt{coneckin.github.io} yang tersinkronisasi langsung dengan API produksi dan rilis APK terbaru.
\end{enumerate}

\section{Rumusan Masalah}
Berdasarkan latar belakang tersebut, rumusan masalah pada Proyek 2 ini adalah sebagai berikut:
\begin{enumerate}[leftmargin=*]
  \item Bagaimana merancang dan melengkapi arsitektur \textit{Hybrid CRUD} secara penuh pada sisi penyimpanan terenkripsi lokal Android serta pada layanan REST API server MongoDB?
  \item Bagaimana menerapkan mekanisme penjaminan mutu berbasis pengujian otomatis (\textit{unit test}, \textit{instrumented test}, dan pengujian integrasi API) dengan ambang \textit{line coverage} minimal 80\% yang terintegrasi dalam pipeline CI/CD?
  \item Bagaimana membangun \textit{landing page} panduan instalasi dan direktori portal berbasis konvensi Croot.js pada \texttt{coneckin.github.io} yang terhubung langsung ke API produksi dan rilis APK terbaru?
  \item Bagaimana mengotomatisasi penyusunan dan rilis dokumentasi akademik LaTeX pada repositori organisasi GitHub Coneckin?
\end{enumerate}

\section{Tujuan Proyek}
Tujuan yang dicapai dalam pelaksanaan Proyek 2 ini meliputi:
\begin{enumerate}[leftmargin=*]
  \item Melengkapi fitur \textit{Ubah Akun} pada aplikasi Android Coneckin serta memisahkan logika bisnis murni agar dapat diuji pada JVM tanpa ketergantungan antarmuka.
  \item Mengimplementasikan endpoint CRUD lengkap pada server Express 5 dan MongoDB (\texttt{coneckin\_db}) yang diamankan otorisasi token admin (SHA-256 \texttt{timingSafeEqual}) serta CORS baca publik.
  \item Mencapai \textit{code coverage} $\ge 80\%$ pada lapisan logika Android (JaCoCo) dan server Node.js, dilengkapi \textit{instrumented test} \texttt{AndroidKeyStore} dan integrasi \textit{test-before-release} di GitHub Actions.
  \item Menayangkan \textit{landing page} responsif berbasis konvensi Croot.js di \url{https://coneckin.github.io/} yang terhubung langsung dengan API produksi di \url{https://panlink.alwaysdata.net}.
  \item Menyediakan alur kompilasi dan rilis otomatis dokumen PDF berbasis LaTeX pada repositori \texttt{laporan} dan \texttt{jurnal}.
\end{enumerate}

\section{Ruang Lingkup dan Struktur Ekosistem}
Ekosistem Coneckin dikelola pada organisasi GitHub \url{https://github.com/coneckin} yang terbagi ke dalam lima repositori resmi sebagaimana dirangkum pada Tabel~\ref{tab:ekosistem}.

\begin{table}[H]
\centering
\caption{Pemetaan Lima Repositori Ekosistem Proyek 2 Coneckin}
\label{tab:ekosistem}
\begin{tabular}{@{}p{3.1cm}p{4.1cm}p{5.6cm}@{}}
\toprule
\textbf{Repositori} & \textbf{Teknologi Utama} & \textbf{Peran dalam Ekosistem} \\
\midrule
\texttt{coneckin} & Kotlin, Jetpack Compose, JUnit 4, JaCoCo; Node.js, Express 5, MongoDB & Aplikasi Android native, API server REST, pengujian otomatis, dan CI/CD rilis APK/AAB. \\
\texttt{coneckin.github.io} & HTML5, CSS3, ES6 Module (konvensi Croot.js) & \textit{Landing page} panduan instalasi, direktori portal dinamis, status API, dan tautan rilis APK terbaru. \\
\texttt{laporan} & LaTeX, GitHub Actions & Laporan akhir Proyek 2 dengan kompilasi dan rilis PDF otomatis. \\
\texttt{jurnal} & LaTeX, GitHub Actions & Artikel ilmiah Proyek 2 dalam format jurnal dengan rilis PDF otomatis. \\
\texttt{presentasi} & Crootss.if.co.id & Dek slide presentasi untuk demo dan sidang Proyek 2. \\
\bottomrule
\end{tabular}
\end{table}

\section{Sistematika Penulisan}
Laporan ini disusun dalam lima bab: Bab~1 membahas latar belakang, rumusan masalah, tujuan, dan ruang lingkup; Bab~2 menguraikan landasan teori dan teknologi yang digunakan; Bab~3 menjelaskan analisis arsitektur \textit{Hybrid CRUD}, keamanan, dan perancangan sistem; Bab~4 memaparkan hasil implementasi, pengujian otomatis, dan verifikasi produksi; serta Bab~5 berisi kesimpulan dan saran pengembangan lanjutan.
